% !TEX root = ../../../main.tex
\chapter{Formal}
\label{cha:form_EN}

The following sections deal with various formal aspects that should essentially be followed. It is essential to adhere to the guidelines on external form, the depth of structure and the citation style.

\section{Language}
\label{sec:language}

Scientific papers at IWES can be written in German or English, whereby there is a separate template or guideline for English-language papers. The choice of language must be agreed with the supervisor of the thesis. In any case, the valid spelling and grammar must be used.

The language style should be appropriate for a scientific paper. Report language is to be used. Colloquial language should be avoided. Scientific writing is characterized by a very efficient language style. It should always be factual and as concise as possible. Unnecessary filler words and transitional formulations should be avoided. Any word that is not absolutely necessary for understanding the statement should be omitted. Of course, care should still be taken to ensure an appealing reading flow.

To minimize the number of errors, tools for effective spelling and grammar checking should be used before submitting the paper. Linguistic mistakes can also lead to a devaluation of the paper.




\section{Outer Form}
\label{sec:form_EN}

Scientific papers should be single-sided, single-spaced and justified with 1.5 line spacing. A sans serif font should be used. The font size must be between 11 pt and 12 pt. Alternatively, if the work is written using \LaTeX{}, the standard font of \LaTeX{} can be used at 11 pt. This is extended to the German character set using the packages \textbackslash{}usepackage[T1]\{fontenc\} and \textbackslash{}usepackage[utf8]\{inputenc\}.

Apart from the cover page and title page, which have their own formatting, the following page margins must be adhered to

\begin{itemize}
	\item 2.5 cm (top, left and right)
	\item 2.0 cm (bottom)
\end{itemize}

Chapters must start with a new page. For sub-chapters and sub-sub-chapters, however, the page break is kept variable.


\section{Outline Levels, Headers and Page Numbering}
\label{sec:numbering}

No more than three numbered outline levels may be used, i.e. chapter, subchapter and sub-subchapter. If a fourth level of structure is required, the paper must be restructured in order to comply with the limit of three levels of structure.

A header must be provided from the title page onwards. The chapter number and chapter title are listed on the left-hand side of the header. The page numbering is to be arranged on the right-hand side of the header. 

The introductory pages (affidavit, assignment, abstract and table of contents) are numbered in Roman numerals (I, II, III, ...). The Roman numeration begins with the title page (first paper page). From the indexes (list of figures etc.) onwards, the Roman numbering must be continued consecutively. 

\begin{addmargin}[1cm]{1cm}
	\textit{\textbf{example:} If the table of contents ends with the Roman page number \textit{VI}, the page numbering in the directories must continue with the Roman page number \textit{VII}.}
\end{addmargin}
	
	The actual main part of the text, i.e. all chapters from the introduction to the summary and the outlook, are numbered with Arabic numerals starting with 1 (1, 2, 3, ...). The numbering starts with the first page of the introduction.



\section{Formulae, Figures and Tables}

The requirements for the use and integration of formulae, images and tables are set out in separate chapters, see the sections \ref{cha:Formulae}, \ref{cha:Images_EN} and \ref{cha:Tables}. Examples of different types of integration are also presented there. Common requirements are defined below.

Formulae, figures and tables must be numbered. Consecutive numbering within the chapter should be used for this purpose. This means that the respective numbers begin with the chapter number and the numbering is consecutive after a subsequent point. The consecutive numbering after the period is independent of each other for formulae, figures and tables and must start at 1 again in each chapter. The principle is summarized in table \ref{tab:numbering}. The outline level within a chapter is not taken into account in the numbering.




\begin{table}[!htbp]
	\small
	\centering
	\caption[Numbering examples for formulae, figures and tables]{Numbering examples for formulae, figures and tables}
	\begin{tabular}{P{2.5cm}P{1.5cm}P{1cm}P{2cm}P{1cm}P{1.5cm}P{1cm}}
		\toprule
		\textbf{chapter} & \textbf{formula} & \textbf{no.} & \textbf{figure} & \textbf{no.} & \textbf{table} & \textbf{No.} \\
		\midrule
		Chapter 1 
		& 1st formula & 1.1 & 1st figure & 1.1 & 1st table & 1.1 \\
		& 2nd formula & 1.2 & 2nd figure & 1.2 & 2nd table & 1.2 \\
		& ... & ... & ... & ... & ... & ... \\
		Chapter 2 
		& 1st formula & 2.1 & 1st figure & 2.1 & 1st table & 2.1 \\
		& 2nd formula & 2.2 & 2nd figure & 2.2 & 2nd table & 2.2 \\
		& ... & ... & ... & ... & ... & ... \\
		Chapter 3
		& 1st formula & 3.1 & 1st figure & 3.1 & 1st table & 3.1 \\
		& 2nd formula & 3.2 & 2nd figure & 3.2 & 2nd table & 3.2 \\
		... & ... & ... & ... & ... & ... & ... \\
		\bottomrule
	\end{tabular}%
	\label{tab:numbering}%
\end{table}%

\section{References and Footnotes}
\label{sec:References}

Within a chapter, references to other parts of the work are often necessary. References should be made using the respective section number. For example, if you want to refer to the content of the section "`Task"', this must be done using the section number, in this case \ref{sec:Task}. 

\begin{addmargin}[1cm]{1cm}
	\textit{\textbf{example:} Information on integrating the task can be found in section \ref{sec:Task}.}
\end{addmargin}

Such references are to be created using the corresponding commands or functions of the word processing programs. Entering the numbers manually is very likely to lead to errors in the final document. 

If you are referring to a page that is not part of a chapter, you can also refer to it by entering the page number. However, this should generally be avoided.

If reference is made to formulae, figures, tables or appendices, this must be done by stating the formula, figure, table or appendix number (see also table \ref{tab:numbering}).

Footnotes\footnote{This is an example of a footnote} are used when a statement is important for the reader's understanding but is superfluous or misplaced in the body text. Footnotes are numbered consecutively throughout the document. The footer is separated by a horizontal line about 6.5 cm long.



\section{Symbols and Abbreviations}
\label{sec:Abbrev}

Symbols (including formula symbols) and abbreviations should be introduced within the body text where they are used for the first time. For symbols, this is discussed again in section \ref{cha:Formulae}. Abbreviations should be inserted in brackets after the full expression. 

\begin{addmargin}[1cm]{1cm}
	\textit{\textbf{example:} This is the guide of the Institute for Wind Energy Systems (IWES) for the preparation of scientific papers}.
\end{addmargin}

\section{Citation Guidelines}
\label{sec:Citation}

In principle, all parts of sentences, sentences or sections taken verbatim from other publications as well as figures or tables taken from other publications must be cited (so-called direct quotations). Quotation marks must be used to indicate parts of the text that have been taken verbatim. In this case, the citation must include the page and line numbers. Statements, illustrations, etc. that have been taken over analogously must also be identified as such by stating the source (indirect quotations). Direct quotations are extremely unusual in engineering and should be avoided. Indirect quotations are generally to be preferred.

Clear references are essential in scientific work. This enables the reader to find them and thus to check the origin of the scientific statements, unless they are independently developed ideas.

A numerical citation style is to be used in scientific papers at IWES. The number of the respective source must be placed in square brackets after the cited statement in the bibliography. 

\begin{addmargin}[1cm]{1cm}
	\textit{\textbf{example:} A detailed description of control concepts for large wind energy systems can be found in \cite{buch}.}
\end{addmargin}

Sources are listed in the bibliography in alphabetical order. It is important that the formatting of entries in the bibliography is consistent throughout. There are different formatting rules for different types of sources. As a general rule, citable sources should be used. In particular, articles in specialist journals \cite{artikel}, specialist books \cite{buch}, articles in conference proceedings \cite{inp}, parts of an anthology \cite{in}, technical documentation \cite{manual}, published reports from a university or other research institution \cite{bericht} and dissertations \cite{thesis} are considered citable. Sources that cannot be cited should always be avoided. Although Internet sources \cite{internet} cannot always be dispensed with, especially for illustrations, original sources (e.g. the sources of Wikipedia articles) should preferably be used. The references cited above are examples. The corresponding formatting within the bibliography should be adhered to.


\begin{addmargin}[1cm]{1cm}
	\textit{\textbf{Note:} Lecture notes are not considered citable sources as they are treated as general knowledge. In case of doubt, the original literature cited in lecture notes should be cited. Internet sources should also be classified as non-citable, as they are constantly changing and may not be verifiable}.
\end{addmargin}


\section{Submission}
\label{sec:submission}

A scientific paper prepared at the IWES must be submitted in printed and bound form in a white cardboard cover with the cover page printed in color. 

\begin{addmargin}[1cm]{1cm}
	\textit{\textbf{Note:} Printed versions in ring binding are not accepted, as an inventory is then not possible}
\end{addmargin}



The digital copy of the thesis must be sent as a PDF file by e-mail to the examiner and the supervisor. Further electronic data, including all files relevant for reproducibility (such as sources, results, tables, Excel calculations, simulation results, etc.), should be submitted via \textit{SeaFile}. The corresponding upload link will be provided by the supervisor on request.

The number of printed versions to be submitted can be found in the assignment.





